\chapter*{Phụ Lục IV. Nhật Ký Tự Răn Ba Mươi Ngày}
\addcontentsline{toc}{chapter}{Phụ Lục IV. Nhật Ký Tự Răn Ba Mươi Ngày}
\markboth{Nhật Ký Tự Răn 30 Ngày}{Nhật Ký Tự Răn 30 Ngày}

\begin{truyen}{Cách Dùng Phụ Lục Này}{How to Use This Appendix}
\chuhoa{M}{ỗi ngày, chỉ cần đọc một dòng nhắc, trả lời một câu hỏi ngắn, rồi ghi lại một việc rất nhỏ mình đã làm để giữ tiền cho đúng.}
\textit{(Each day, simply read one reminder line, answer one short question, and write down one very small act you did to keep money wisely.)}

Không cần viết dài. Điều quan trọng là sự đều đặn. Ba mươi ngày đủ để một ý nghĩ bắt đầu thành nếp.\
\textit{(There is no need to write long entries. What matters is consistency. Thirty days are enough for an idea to begin becoming a habit.)}
\end{truyen}

\section*{Ngày 1 đến Ngày 10}
\begin{truyen}{Mười Ngày Đầu: Học Nhìn Thẳng}{First Ten Days: Learn to See Clearly}
Ngày 1: Hôm nay tôi đã tiêu gì vì cảm xúc?\
\textit{(Day 1: What did I spend on today because of emotion?)}

Ngày 2: Khoản nào nhỏ nhưng lặp lại đang làm tôi mất tiền nhiều nhất?\
\textit{(Day 2: Which small repeating expense is draining me most?)}

Ngày 3: Tôi còn nhớ đồng tiền của mình đến từ đâu không?\
\textit{(Day 3: Do I still remember where my money comes from?)}

Ngày 4: Tôi có đang tiêu để quên mệt không?\
\textit{(Day 4: Am I spending in order to forget my tiredness?)}

Ngày 5: Điều gì hôm nay là cần thật, điều gì chỉ là thích?\
\textit{(Day 5: What today was truly needed, and what was merely wanted?)}

Ngày 6: Tôi có đang nhầm sĩ diện với nhu cầu không?\
\textit{(Day 6: Am I confusing pride with necessity?)}

Ngày 7: Tôi đã ghi chép đủ thành thật chưa?\
\textit{(Day 7: Have I recorded my spending honestly enough?)}

Ngày 8: Tôi thường yếu lòng nhất vào thời điểm nào trong ngày?\
\textit{(Day 8: At what time of day am I most vulnerable?)}

Ngày 9: Điều gì làm tôi thấy thiếu dù tôi chưa thật sự thiếu?\
\textit{(Day 9: What makes me feel lacking even when I am not truly lacking?)}

Ngày 10: Một món tôi chưa mua hôm nay đã cứu tôi khoản nào?\
\textit{(Day 10: What did one unbought item save me today?)}
\end{truyen}

\section*{Ngày 11 đến Ngày 20}
\begin{truyen}{Mười Ngày Giữa: Học Giữ Tay}{Middle Ten Days: Learn to Hold the Hand}
Ngày 11: Hôm nay tôi đã dừng lại trước món gì?\
\textit{(Day 11: What did I stop myself from buying today?)}

Ngày 12: Tôi đã để dành được bao nhiêu, dù rất ít?\
\textit{(Day 12: How much did I manage to save, even if very little?)}

Ngày 13: Tôi có đang đợi còn dư mới tiết kiệm không?\
\textit{(Day 13: Am I still waiting for leftovers before saving?)}

Ngày 14: Hôm nay tôi có tiêu gì chỉ để tự an ủi?\
\textit{(Day 14: Did I buy anything today only to comfort myself?)}

Ngày 15: Nếu món này phải đổi bằng năm giờ đứng lớp, tôi còn muốn nó không?\
\textit{(Day 15: If this item costs five hours of teaching, do I still want it?)}

Ngày 16: Quỹ dự phòng của tôi đang ở đâu?\
\textit{(Day 16: Where does my emergency fund stand today?)}

Ngày 17: Tôi đã nói chuyện tiền bạc rõ ràng với gia đình chưa?\
\textit{(Day 17: Have I spoken clearly with my family about money?)}

Ngày 18: Tôi có khoản vay nào đang sinh ra từ thói quen xấu không?\
\textit{(Day 18: Do I have any debt created by a bad habit?)}

Ngày 19: Hôm nay tôi đã dạy con điều gì bằng hành động của mình?\
\textit{(Day 19: What did I teach my child today through my actions?)}

Ngày 20: Tôi đã biết ơn điều gì thay vì đòi thêm?\
\textit{(Day 20: What did I feel grateful for instead of demanding more?)}
\end{truyen}

\section*{Ngày 21 đến Ngày 30}
\begin{truyen}{Mười Ngày Cuối: Học Xây Nếp Mới}{Final Ten Days: Learn to Build a New Pattern}
Ngày 21: Một lỗ rò nào tôi đã bịt lại được?\
\textit{(Day 21: Which leak have I managed to close?)}

Ngày 22: Một thói quen nào tôi vẫn còn đang nuông chiều?\
\textit{(Day 22: Which habit am I still indulging?)}

Ngày 23: Hôm nay tôi sống vừa sức ở điểm nào?\
\textit{(Day 23: Where did I live within my means today?)}

Ngày 24: Tôi có đang giữ ngày mai bằng một khoản nhỏ hôm nay không?\
\textit{(Day 24: Am I protecting tomorrow through a small act today?)}

Ngày 25: Điều gì từng làm tôi tiêu tiền mù quáng nhưng nay đã bớt sức kéo?\
\textit{(Day 25: What used to make me spend blindly but now has less pull?)}

Ngày 26: Tôi có đang mệt nên muốn mua sắm cho quên không?\
\textit{(Day 26: Am I tired and wanting to shop just to forget?)}

Ngày 27: Hôm nay tôi đã chọn bình yên thay vì hình ảnh ở điểm nào?\
\textit{(Day 27: Where did I choose peace over image today?)}

Ngày 28: Nếu tháng này kết thúc ngay hôm nay, tôi có hối tiếc cách tiêu của mình không?\
\textit{(Day 28: If this month ended today, would I regret the way I spent?)}

Ngày 29: Tôi muốn con mình học điều gì từ ví tiền của cha mẹ?\
\textit{(Day 29: What do I want my children to learn from their parents’ wallet?)}

Ngày 30: Tôi sẽ giữ lại thói quen tốt nào cho ba mươi ngày tiếp theo?\
\textit{(Day 30: Which good habit will I carry into the next thirty days?)}
\end{truyen}

\begin{baihoc}
Thói quen không hình thành trong một ngày. Nhưng ba mươi ngày thành thật với chính mình đủ để một con đường mới bắt đầu rõ ra.
\textit{(Habits do not form in a single day. But thirty days of honesty with oneself are enough for a new path to begin showing itself clearly.)}
\end{baihoc}

\begin{ghinhoanh}
Thirty honest days can redirect a life.

Small daily truth becomes discipline.
\end{ghinhoanh}

\ngancach